\section{A complex associated to simple normal crossing divisors}\label{complex_for_SNC}
We now discuss a complex that, as we will see later, gives
a filtered resolution of $\jepPubScript{M}_r(h^{-\alpha})$ by filtered induced $\jepPubScript{D}_X$-modules in the case when $h$ defines a simple normal crossing divisor. 


Let $X$ be a smooth, $n$-dimensional, complex variety, $h\in\jepPubScript{O}_X(X)$ nonzero, and $\alpha$ a nonzero rational number (we allow $\alpha$ 
to be either positive or negative). Let $D = \alpha \cdot {\rm div}(h)$. We denote by $Z$
the support of $D$, and assume that it has simple normal crossings.


Associated to $Z$ we have the following complex of right $\jepPubScript{D}_X$-modules:
$$C^{\jepPubBullet}:\,\,0\longrightarrow \jepPubScript{D}_X\longrightarrow\Omega_X^1(\log Z)\otimes_{\jepPubScript{O}_X}\jepPubScript{D}_X\longrightarrow\cdots\longrightarrow \Omega_X^n(\log Z)\otimes_{\jepPubScript{O}_X}\jepPubScript{D}_X\longrightarrow 0,$$
placed in degrees $-n,\ldots,0$. We denote by $D_i\colon C^i\to C^{i+1}$ its differentials. 
If $x_1,\ldots,x_n$ are local coordinates on $X$, then
$$D_i(\eta\otimes P)=d\eta\otimes P+\sum_{i=1}^n (dx_i\wedge\eta)\otimes \partial_{x_i}P.$$
In fact $C^{\jepPubBullet}$ is a filtered complex, where
$$F_{p-n}C^i=\Omega_X^{i+n}(\log Z)\otimes_{\jepPubScript{O}_X}F_{p+i}\jepPubScript{D}_X.$$
This filtered complex is quasi-isomorphic to the filtered right $\jepPubScript{D}_X$-module $\omega_X(*Z)$ 
corresponding to the filtered left $\jepPubScript{D}_X$-module $\jepPubScript{O}_X(*Z)$ (see \cite[Prop.~3.1]{MP1}, and 
\cite[Prop.~3.11(ii)]{Saito-MHM} for a more general statement).

Given $h$ and $\alpha$ as above, we also consider the filtered complex $C^{\jepPubBullet}_{h^{-\alpha}}$ consisting of the same sheaves, but 
with differential $C^i_{h^{-\alpha}}\to C^{i+1}_{h^{-\alpha}}$ given by
$$D_i-\big((\alpha\cdot {\rm d\log}(h)\wedge \jepPubBullet)\otimes {1_{\jepPubScript{D}_X}}\big).\footnote
{In related settings, for instance involving the de Rham complex of $\jepPubScript{M}(h^{-\alpha})$, this type of complex can already be found in the literature; see for instance \cite[\S6.3.11]{Bjork}.}$$
It is easy to see that this is indeed a filtered complex. 

Suppose now that we also have an effective divisor $T$ supported on $Z$. It is not hard to check that the formula for the map
$$C^i_{h^{-\alpha}}\longrightarrow C^{i+1}_{h^{-\alpha}}$$
induces also a map
$$C^i_{h^{-\alpha}}(-T):=\jepPubScript{O}_X(-T)\otimes_{\jepPubScript{O}_X}\Omega^{i+n}_X(\log Z)\otimes_{\jepPubScript{O}_X}\jepPubScript{D}_X$$
$$\longrightarrow 
C^{i+1}_{h^{-\alpha}}(-T):=\jepPubScript{O}_X(-T)\otimes_{\jepPubScript{O}_X}\Omega^{i+1+n}_X(\log Z)\otimes_{\jepPubScript{O}_X}\jepPubScript{D}_X.$$
This is due to the fact that if locally $T={\rm div}(u)$ and $\eta$ is a local section of $\Omega^{i+n}_X(\log Z)$,
then we can write 
$d(u\eta)=ud(\eta)+u\cdot {\rm d\log}(u)\wedge\eta$.
We thus obtain a filtered subcomplex $C^{\jepPubBullet}_{h^{-\alpha}}(-T)$ of $C^{\jepPubBullet}_{h^{-\alpha}}$.
We emphasize that this is \emph{not} obtained by tensoring $C^{\jepPubBullet}_{h^{-\alpha}}$ with $\jepPubScript{O}_X(-T)$.


\begin{proposition}\label{prop_SNC_complex}
If no coefficient of $D - T$ lies in $\jepPubBold{Z}_{<0}$, then
the complex $C^{\jepPubBullet}_{h^{-\alpha}}(-T)$
is filtered quasi-isomorphic to $\big(h^{-\alpha}\omega_X(*Z), G_\jepPubBullet \big)$, where 
$$G_{k-n}h^{-\alpha}\omega_X(*Z)=0\quad\text{if}\quad k<0,$$
$$G_{-n}h^{-\alpha}\omega_X(*Z)=h^{-\alpha}\omega_X(Z- T)$$
$$\text{and}\quad G_{k-n}h^{-\alpha}\omega_X(*Z)=G_{-n}h^{-\alpha}\omega_X(*Z)\cdot F_k\jepPubScript{D}_X\quad\text{if}\quad k>0.$$
\end{proposition}

\begin{proof}
It is immediate to check that the differential induced on ${\rm gr}^F_pC^{\jepPubBullet}_{h^{-\alpha}}(-T)$ does become 
equal to the differential $D_i$ twisted with the identity on $\jepPubScript{O}_X(-T)$, and therefore for every $p$ we have
$${\rm gr}^F_pC^{\jepPubBullet}_{h^{-\alpha}}(-T)=\jepPubScript{O}_X(-T)\otimes_{\jepPubScript{O}_X}{\rm gr}^F_pC^{\jepPubBullet}.$$
In particular, we have 
$$H^iF_pC^{\jepPubBullet}_{h^{-\alpha}}(-T)=0\quad\text{for every}\quad p\in\jepPubBold{Z}\,\,\text{and}\,\, i\in\jepPubBold{Z}\smallsetminus\{0\},$$
by the result in \cite{MP1} quoted above. Consider now the morphism of right $\jepPubScript{D}_X$-modules
$$\varphi\colon C^0_{h^{-\alpha}}(-T) = \omega_X (Z - T)\otimes_{\jepPubScript{O}_X} \jepPubScript{D}_X \longrightarrow 
h^{-\alpha}\omega_X(*Z)$$
given by 
$$\varphi(w\otimes\eta\otimes Q)=(h^{-\alpha}w\eta)Q.$$
We first check that this morphism is surjective. We do this locally, hence we may assume that we have a system of coordinates $x_1,\ldots,x_n$ on $X$ such that $\jepPubScript{O}_X(-Z)$ is generated by $x_1\cdots x_r$ and $\jepPubScript{O}_X(-T)$ by $x_1^{\beta_1}\cdots
x_r^{\beta_r}$. We also write $h=u x_1^{a_1}\cdots x_r^{a_r},$ where $u$ is an everywhere nonvanishing function, 
and define $\alpha_i=\alpha a_i$ and $\gamma_i=\alpha_i-\beta_i$ for all $i$. Note for later use that
$$\alpha\cdot {\rm d\log}(h)=\frac{\alpha}{u}du+\sum_{i=1}^r\alpha_i\frac{dx_i}{x_i}.$$

The surjectivity of $\varphi$ follows from the fact that 
\begin{equation}\label{eq0_prop_SNC_complex}
{\rm Im}(\varphi)=(h^{-\alpha}x_1^{\beta_1}\cdots x_r^{\beta_r}\eta)\cdot\jepPubScript{D}_X=h^{-\alpha}\omega_X(*Z),
\end{equation}
where $$\eta={\rm d\log}(x_1)\wedge\cdots\wedge
{\rm d\log}(x_r)\wedge dx_{r+1}\wedge\cdots\wedge dx_n,$$ and the second equality in (\ref{eq0_prop_SNC_complex})
   is a consequence of the fact that $-\gamma_i-1\not\in \jepPubBold{Z}_{\geqslant 0}$ for all $i$, by assumption.

In order to complete the proof of the proposition it is enough to show
that, for every $k\geqslant 0$, the following sequence is exact:
$$
\jepPubScript{O}_X(-T)\otimes\Omega_X^{n-1}(\log Z)\otimes F_{k-1}\jepPubScript{D}_X\xrightarrow{\;\textstyle\psi_k\;}\jepPubScript{O}_X(-T)\otimes \omega_X(Z)\otimes F_k\jepPubScript{D}_X$$ 
$$\xrightarrow{\;\textstyle\varphi_k\;} G_{k-n} h^{-\alpha}\omega_X(*Z) \longrightarrow 0,
$$
where $\varphi_k$ is the restriction of $\varphi$ to the $(k-n)$-th level of the filtration and $\psi_k$ is the restriction of the  differential of  $C^{\jepPubBullet}_{h^{-\alpha}}(-T)$.
The surjectivity of $\varphi_k$ is an immediate consequence of the surjectivity of $\varphi$ and the definition of the filtration on $h^{\alpha}\omega_X(*Z)$.



Keeping the above notation for the local coordinates on $X$,
it follows from the definition of $\psi_k$ that
$${\rm Im}(\psi_k)=\prod_{j=1}^rx_j^{\beta_j}\otimes \eta \otimes \bigg(\sum_{i=1}^r\Big(x_i\partial_i-\gamma_i-\frac{\partial u}{\partial x_i}\cdot \frac{\alpha x_i}{u}\Big)\cdot F_{k-1}\jepPubScript{D}_X $$
$$ + \sum_{i=r+1}^n
\Big(\partial_i-\frac{\partial u}{\partial x_i}\cdot\frac{\alpha}{u}\Big)\cdot F_{k-1}\jepPubScript{D}_X\bigg)$$
and it is straightforward to see that this is contained in ${\rm Ker}(\varphi_k)$. We now prove by induction on $k$
that if $\varphi_k(x_1^{\beta_1}\cdots x_r^{\beta_r}\otimes\eta\otimes P)=0$ for some $P\in F_k\jepPubScript{D}_X$, then $x_1^{\beta_1}\cdots x_r^{\beta_r}\otimes\eta \otimes P\in {\rm Im}(\psi_k)$. 
Note that the case $k=0$ is trivial.
Let us write $P=\sum_{u,v}c_{u,v}\partial^{u}x^{v}$, where $u$ and $v$ vary over $\jepPubBold{Z}_{\geqslant 0}^n$. 
After subtracting suitable terms from $P$, we may assume that whenever $c_{u,v}\neq 0$, we have $u_i=0$ for $i>r$.
Furthermore, note that if $u_i,v_i>0$ for some $i\leqslant r$, then we can write
$$\partial^{u}x^{v}=\Big(x_i\partial_i-\gamma_i-\frac{\partial u}{\partial x_i}\cdot \frac{\alpha x_i}{u}\Big)A+B,$$
with both $A$ and $B$ of order $\leqslant k-1$. Therefore we may also assume that whenever $c_{u,v}\neq 0$ and $|u|:=\sum_iu_i=k$, we have
\begin{equation}\label{eq_prop_SNC_complex}
u_iv_i=0\quad\text{for $i$ such that}\quad 1\leqslant i\leqslant n.
\end{equation}



Since 
$$(h^{-\alpha}x_1^{\beta_1}\cdots x_r^{\beta_r}\eta)\partial^{u}x^{v}=(\text{non-zero constant}\footnote{Here we use again the fact that
$-\gamma_i-1\not\in\jepPubBold{Z}_{\geqslant 0}$ for all $i$.})\cdot (h^{-\alpha}x_1^{\beta_1}\cdots x_r^{\beta_r}\eta)x^{v-u},$$
and since (\ref{eq_prop_SNC_complex}) implies that for every $(u,v)$ and $(u',v')$ with $|u|=k$ and $c_{u,v},c_{u',v'}\neq 0$
we have $x^{v-u}\neq x^{v'-u'}$, we conclude that in fact $P\in F_{k-1}\jepPubScript{D}_X$, hence we are done by induction.
\end{proof}










